\section{Introduction}
\label{sec:intro}

Upcoming exascale clusters will be evaluated on large-scale LLM pre-training. The work started as the
validation of two new milabench workloads at a scale they had never run at. Full access to Rorqual extends
it: the same workloads serve as probes of how the cluster, and its network architecture in particular,
behaves under large-scale training, which is what its administrators and the exascale evaluation need to
know. Rorqual is compared with Fir and Tamia, two clusters with the same GPU nodes and
four times the inter-node bandwidth per node (\cref{tab:clusters}).

\begin{table}[h]
  \caption{Clusters. All nodes have 4 H100 80\,GB GPUs connected by NVLink.}
  \label{tab:clusters}
  \centering\small
  \begin{adjustbox}{max width=\linewidth}
  \begin{tabular}{llll}
    \toprule
    Cluster & Nodes used & Inter-node network per node & Role \\
    \midrule
    Rorqual & 66 (264 GPUs), 16 leaf switches & $1\times$ 200\,Gb/s HDR InfiniBand (50\,Gb/s per GPU) & target, full-access window \\
    Fir     & leaves of 20 nodes & $4\times$ 200\,Gb/s (one per GPU) & preliminary runs \\
    Tamia   & 5 leaves of 8--16 nodes & $4\times$ 200\,Gb/s & packed vs.\ spread placement \\
    \bottomrule
  \end{tabular}
  \end{adjustbox}
\end{table}

\paragraph{Rorqual topology.} Two levels: 25 leaf switches under one spine switch. Of the 93 GPU nodes, 27
are MIG-partitioned (9 leaves) and unusable here, leaving 66 full nodes: 13 leaves with 3 GPU nodes each,
and three larger leaves with 7, 8 and 12 (ihpc218, ihpc217, ihpc316). Every leaf except ihpc316 also
connects CPU nodes, whose jobs share its uplinks.

\paragraph{Workloads.}
\begin{itemize}[leftmargin=*, itemsep=1pt]
  \item \textbf{Qwen3.8-27B} (dense, 64 layers: 48 Gated DeltaNet + 16 gated attention), TorchTitan,
        FSDP over all GPUs, $3 \times 4{,}096$ tokens per GPU per step (weak scaling), BF16 compute with
        FP32 master weights and AdamW moments, full activation checkpointing, regional compilation.
  \item \textbf{GLM-5} (mixture of experts), Megatron via slime: full model at about 208--264 GPUs,
        truncated versions at smaller scale. \placeholder{configuration once the pipeline is fixed.}
  \item \textbf{Communication benchmarks} (nccl-tests all-reduce and all-to-all) and \textbf{health
        checks} before each allocation.
\end{itemize}

\paragraph{Questions.} Each section answers one question:
(\cref{sec:hardware}) do nodes perform as expected and uniformly?
(\cref{sec:fabric}) how much does the network limit training, and does placement on leaf switches matter?
(\cref{sec:scaling}) how does throughput per GPU change from one switch to 264 GPUs?
(\cref{sec:moe}) how does an MoE model behave at scale?
(\cref{sec:software}) how does the software stack compare with expectations and published numbers?
(\cref{sec:operations}) is the system ready for production runs (startup, failures, variance, stability)?

\paragraph{Declared deviations} from the reference training recipe, valid for every result unless noted:
AdamW instead of Muon; multi-token prediction (MTP) off; random tokens instead of a corpus; FP8 not
used for Qwen (not supported on our TorchTitan pin; used for GLM-5).
